
\section{Introduction}
\label{sec:intro}

Since the introduction of Sora~\citep{openaisora2024} by OpenAI, video generation technology has attracted substantial attention from both industry and academia, leading to rapid advancements in the field.
%
The emergence of models capable of generating videos that are on par with professionally crafted content has significantly improved the efficiency of content creation while simultaneously reducing the costs associated with video production.
%
These rapid advancements in video generation technology have also been greatly attributed to the development of the open-source community.
%
Notable projects like HunyuanVideo~\citep{kong2024hunyuanvideo}, Mochi~\citep{genmo2024mochi}, and CogVideoX~\citep{cogvideox} have made their video foundation model codes and weights publicly available, gradually narrowing the gap between open-source models and commercial counterparts.
%
However, it is essential to acknowledge the persistent gap between these outstanding open-source models and the latest closed-source models. 
%
This gap is primarily evident in three aspects:
%
\emph{Suboptimal Performance}: A notable performance gap remains, as the pace of development in commercial models far exceeds that of current open-source models, resulting in significantly superior capabilities.
%
\emph{Limited Capabilities}: Most foundational models are limited to general text-to-video (T2V) tasks, whereas the demands of video creation are multifaceted. Consequently, basic T2V models are insufficient to address these diverse requirements.
%
\emph{Insufficient Efficiency}: Despite their impressive performance and scale, these models often prove impractical for creative teams with limited computational resources, hindering their accessibility and usability.
%
These challenges collectively impose constraints on the continued growth and innovation within the open-source community.


To address the aforementioned challenges, this report introduces and publicly releases a novel series of high-performance foundational video generation models, referred to as \method, which sets a new benchmark in the field.
%
The core design of \method is inspired by the success of Diffusion Transformers (DiT)~\citep{dit} combined with Flow Matching~\citep{flow-matching}, a framework that has demonstrated the significant performance gains achievable through scaling in both text-to-image (T2I)~\citep{esser2024sd3} and text-to-video (T2V)~\citep{kong2024hunyuanvideo} tasks.
%
Within this architectural paradigm, cross-attention is employed to embed text conditions, while the model's design is meticulously optimized to ensure computational efficiency and precise text controllability. 
%
To further enhance the model's ability to capture complex dynamics, a full spatio-temporal attention mechanism is incorporated.
%
Through extensive experimentation, the model is validated at scale, reaching 14 billion parameters.
%
Subsequently, \method has seen large-scale data comprising billions of images and videos, amounting to $\mathcal{O}(1)$ trillions of tokens in total.
%
This extensive training facilitates the emergence of the model's capabilities, allowing it to demonstrate robust performance across multiple dimensions, such as motion amplitude and quality, visual text generation, camera control, instruction adherence, and stylistic diversity.
%
Building upon the powerful foundational model, we have expanded its capabilities to numerous downstream tasks, including image-to-video generation (I2V), instruction-guided video editing (V2V), zero-shot personalized customization, real-time video generation, and audio generation, among other critical applications.
%
To minimize inference costs, we also introduce a 1.3B model alongside a 14B model for T2V and I2V, both of which support 480p resolution and greatly enhance inference efficiency. 
%
Remarkably, the 1.3B model requires only 8.19G of VRAM, allowing it to run on many consumer-grade GPUs, while its performance exceeds that of many larger open-source models.


Moreover, we will also publicly present the entire training process, including the large-scale data construction pipeline, video variational autoencoder (VAE), training strategies, acceleration techniques, and automated evaluation algorithms, to empower the community in developing specialized foundational video models.
%
In addition, we will provide comprehensive design details and experimental results, offering insights into phenomena observed during the resource-intensive training of large generative models alongside our key findings and conclusions.
%
We are confident that these contributions will play a pivotal role in accelerating the advancement of video generation technology.


\begin{figure}[t]  
    \centering  
    \includegraphics[width=1.0\textwidth]{figures/overall/fig02_example.pdf}  
    \caption{Samples generated by \method. \method can produce videos with large motions, high fidelity, and realistic details. It is also the first model to generate both Chinese and English text within videos. Additionally, \method offers a range of features, including text-to-video, image-to-video, and video editing capabilities.} 
    \vspace{-2mm}
    \label{fig:wan_examples_01} % 
\end{figure}


